\section*{Problem 1}

在多元线性回归中，假设设计矩阵 $X$ 列满秩。为了极小化最小二乘目标函数
\[
    \left\|X\bm{\beta} - \bm{y}\right\|^2,
\]
请推导参数 $\bm{\beta}$ 的最优解所满足的正规方程（Normal Equation）。

\begin{proof}
    记最小二乘目标函数为
    \[
        f(\bm{\beta})
        = \left\|X\bm{\beta} - \bm{y}\right\|^2
        = \left(X\bm{\beta} - \bm{y}\right)^\top
          \left(X\bm{\beta} - \bm{y}\right).
    \]
    展开并对 $\bm{\beta}$ 求梯度，得到
    \[
        \begin{aligned}
            f(\bm{\beta})
            &= \bm{\beta}^\top X^\top X\bm{\beta}
               - 2\bm{y}^\top X\bm{\beta}
               + \bm{y}^\top\bm{y}, \\
            \nabla_{\bm{\beta}} f
            &= 2X^\top X\bm{\beta} - 2X^\top\bm{y}
             = 2X^\top\left(X\bm{\beta} - \bm{y}\right).
        \end{aligned}
    \]
    因此最优解 $\bm{\beta}^*$ 满足
    \[
        \nabla_{\bm{\beta}}f(\bm{\beta}^*) = \bm{0}
        \iff
        \boxed{X^\top X\bm{\beta}^* = X^\top\bm{y}},
    \]
    这就是正规方程。

    由于 $X$ 列满秩，对任意非零向量 $\bm{v}$ 都有
    \[
        \bm{v}^\top X^\top X\bm{v} = \|X\bm{v}\|^2 > 0,
    \]
    故 $X^\top X$ 正定且可逆。于是目标函数存在唯一的全局最优解
    \[
        \boxed{\bm{\beta}^* = \left(X^\top X\right)^{-1}X^\top\bm{y}}.
    \]
\end{proof}

\section*{Problem 2}

在单层神经网络
\[
    \bm{z} = \sigma\left(W\bm{x} + \bm{b}\right)
\]
中，设激活函数为 Sigmoid 函数
\[
    \sigma(t) = \frac{1}{1 + e^{-t}}.
\]
请推导导数 $\sigma'(t)$，并将其恒等变形为仅包含 $\sigma(t)$ 的表达式。

\begin{proof}
    对 $\sigma(t)$ 直接求导，得到
    \[
        \sigma'(t)
        = -\left(1 + e^{-t}\right)^{-2}\left(-e^{-t}\right)
        = \frac{e^{-t}}{\left(1 + e^{-t}\right)^2}.
    \]
    另一方面，
    \[
        1 - \sigma(t)
        = 1 - \frac{1}{1 + e^{-t}}
        = \frac{e^{-t}}{1 + e^{-t}}.
    \]
    因此
    \[
        \boxed{\sigma'(t)
        = \frac{1}{1 + e^{-t}}\frac{e^{-t}}{1 + e^{-t}}
        = \sigma(t)\left(1 - \sigma(t)\right)}.
    \]
\end{proof}

\section*{Problem 3}

在主成分分析（PCA）中，求解半正定矩阵 $A^\top A$ 的最大特征值问题等价于求解约束优化问题
\[
    \max_{\|\bm{x}\| = 1}\bm{x}^\top A^\top A\bm{x}.
\]
请利用拉格朗日（Lagrange）乘子法证明：该优化问题的最大值即为矩阵 $A^\top A$ 的最大特征值 $\lambda_1$。

\begin{proof}
    记 $B = A^\top A$。由于 $B$ 是实对称半正定矩阵，其特征值均为非负实数。引入拉格朗日乘子 $\lambda$，构造
    \[
        L(\bm{x}, \lambda)
        = \bm{x}^\top B\bm{x}
          - \lambda\left(\bm{x}^\top\bm{x} - 1\right).
    \]
    对 $\bm{x}$ 求梯度并令其为零，得到
    \[
        \nabla_{\bm{x}}L
        = 2B\bm{x} - 2\lambda\bm{x}
        = \bm{0}
        \iff
        B\bm{x} = \lambda\bm{x}.
    \]
    单位球面是紧集，目标函数连续，故最大值必然取得。因此每个约束极值点 $\bm{x}$ 都是 $B$ 的单位特征向量，对应的拉格朗日乘子 $\lambda$ 是 $B$ 的特征值。又因为 $\|\bm{x}\| = 1$，
    \[
        \bm{x}^\top B\bm{x}
        = \bm{x}^\top\left(\lambda\bm{x}\right)
        = \lambda\bm{x}^\top\bm{x}
        = \lambda.
    \]
    故目标函数在各约束极值点的值恰为 $B$ 的相应特征值。取其中最大的特征值 $\lambda_1$ 及其单位特征向量 $\bm{v}_1$，便有
    \[
        \max_{\|\bm{x}\| = 1}\bm{x}^\top A^\top A\bm{x}
        = \bm{v}_1^\top A^\top A\bm{v}_1
        = \boxed{\lambda_1}.
    \]
\end{proof}
